\begin{abstract}
La integración de Modelos de Lenguaje de Gran Escala (LLM) y arquitecturas de Generación Aumentada por Recuperación (RAG) ha cambiado drásticamente el desarrollo de software. No obstante, la presión por lanzar estas herramientas rápidamente ha llevado a los desarrolladores a ignorar principios básicos de diseño, generando una nueva forma de deuda técnica. Mezclar componentes de inteligencia artificial con código tradicional suele causar fallos estructurales graves; por ejemplo, la dependencia extrema de un solo proveedor, la exposición de las bases de datos vectoriales y la falta de validación de los datos que ingresan al modelo. Para abordar esto, el artículo revisa el estado del arte para caracterizar estos nuevos antipatrones arquitectónicos en aplicaciones GenAI. Al analizar la literatura, se hace evidente que, aunque ya existen estudios que catalogan estos errores, hoy en día solo se detectan de manera manual o cuando el sistema ya falló. A partir de este vacío, el documento justifica la necesidad urgente de contar con soluciones proactivas y sienta las bases teóricas para construir una herramienta de análisis estático. Este motor, apoyado en Árboles de Sintaxis Abstracta (AST) y Minería de Repositorios (MSR), permitirá escanear el código y corregir esta degradación desde las primeras etapas de desarrollo.
\end{abstract}

\begin{IEEEkeywords}
Arquitectura de Software, Antipatrones, GenAI, RAG, Análisis Estático, Deuda Técnica, Mining Software Repositories.
\end{IEEEkeywords}